% GENERATED FILE. Edit canonical lessons, then run npm run generate:books.
% canonical-source-hash: fnv1a64:b3886b4fcc613947
% canonical-lessons: ES-C681-mercadillo, ES-C681-centro-comercial, ES-C681-grandes-almacenes, ES-C681-festivo, ES-C681-atender

\chapter{De compras}
\label{ch:es-de-compras}
\hlchaptermodality{681}

\begin{chapteropening}
\textbf{By the end of this chapter:} \emph{I can shop around a town: the street market, the shopping centre, the department store, holiday opening and who serves me.}

Complete the last lesson of chapter 681: I can shop around a town: the street market, the shopping centre, the department store, holiday opening and who serves me.
\end{chapteropening}

\section[el mercadillo]{el mercadillo — the street market, the flea market}

\label{lesson:ES-C681-mercadillo}



\begin{quote}
Before the new one: say the Spanish for the launderette, then the Spanish for the ice-cream shop.
\end{quote}

\subsection*{You'll want to know: el mercadillo}
\textbf{el mercadillo} — \textquotedblleft{}the street market, the flea market\textquotedblright{}.

\begin{quote}
Los domingos hay mercadillo en la plaza.
\end{quote}

\begin{quote}
\emph{On Sundays there is a street market in the square.}
\end{quote}

Smaller than \emph{el mercado}, often outdoors and once a week.

\begin{cousinweb}
\emph{mercadillo} is a little \emph{mercado} \textquotedblleft{}market\textquotedblright{}, which is inherited from Latin \emph{mercātus} (EtymDB).
\end{cousinweb}

\begin{grammarlens}[title={What you've built}]
One more word for telling, explaining and describing at B1.
\end{grammarlens}

\subsection*{Your turn}
\begin{itemize}
\raggedright
  \item \emph{Say it:} \emph{el mercadillo}
  \item \emph{Say it:} \emph{el mercadillo}, once more
  \item \emph{Say it:} a sentence of your own with \emph{el mercadillo}
  \item \emph{Recall:} say the Spanish for the launderette, then the Spanish for the ice-cream shop, then say \emph{el mercadillo} again
\end{itemize}

\begin{tcolorbox}[breakable,colback=teal!4,colframe=teal!35!black,title={Before you move on}]
What does el mercadillo mean? (\textquotedblleft{}The street market, the flea market\textquotedblright{}.) Say it once more.
\end{tcolorbox}
\section[el centro comercial]{el centro comercial — the shopping centre, the mall}

\label{lesson:ES-C681-centro-comercial}



\begin{quote}
Before the new one: say the Spanish for the ice-cream shop, then the Spanish for the street market.
\end{quote}

\subsection*{You'll want to know: el centro comercial}
\textbf{el centro comercial} — \textquotedblleft{}the shopping centre, the mall\textquotedblright{}.

\begin{quote}
El centro comercial abre hasta las diez.
\end{quote}

\begin{quote}
\emph{The shopping centre is open until ten.}
\end{quote}

In parts of Latin America also \emph{el mall} or \emph{el shopping}.

\begin{cousinweb}
\emph{centro} comes from Latin \emph{centrum}, Greek \emph{kéntron} \textquotedblleft{}a sharp point\textquotedblright{}; \emph{comercial} is built on \emph{comercio}, a learned borrowing from Latin \emph{commercium} \textquotedblleft{}trade\textquotedblright{}.
\end{cousinweb}

\begin{grammarlens}[title={What you've built}]
One more word for telling, explaining and describing at B1.
\end{grammarlens}

\subsection*{Your turn}
\begin{itemize}
\raggedright
  \item \emph{Say it:} \emph{el centro comercial}
  \item \emph{Say it:} \emph{el centro comercial}, once more
  \item \emph{Say it:} a sentence of your own with \emph{el centro comercial}
  \item \emph{Recall:} say the Spanish for the ice-cream shop, then the Spanish for the street market, then say \emph{el centro comercial} again
\end{itemize}

\begin{tcolorbox}[breakable,colback=teal!4,colframe=teal!35!black,title={Before you move on}]
What does el centro comercial mean? (\textquotedblleft{}The shopping centre, the mall\textquotedblright{}.) Say it once more.
\end{tcolorbox}
\section[los grandes almacenes]{los grandes almacenes — the department store}

\label{lesson:ES-C681-grandes-almacenes}



\begin{quote}
Before the new one: say the Spanish for the street market, then the Spanish for the shopping centre.
\end{quote}

\subsection*{You'll want to know: los grandes almacenes}
\textbf{los grandes almacenes} — \textquotedblleft{}the department store\textquotedblright{}.

\begin{quote}
Compré la maleta en unos grandes almacenes.
\end{quote}

\begin{quote}
\emph{I bought the suitcase in a department store.}
\end{quote}

Always plural, even for one store.

\begin{cousinweb}
\emph{almacén} comes from Andalusian Arabic \emph{al-maxzán} \textquotedblleft{}the storehouse\textquotedblright{}, from Arabic \emph{makhzan}; \emph{grande} is inherited from Latin \emph{grandis} \textquotedblleft{}big\textquotedblright{}.
\end{cousinweb}

\begin{grammarlens}[title={What you've built}]
One more word for telling, explaining and describing at B1.
\end{grammarlens}

\subsection*{Your turn}
\begin{itemize}
\raggedright
  \item \emph{Say it:} \emph{los grandes almacenes}
  \item \emph{Say it:} \emph{los grandes almacenes}, once more
  \item \emph{Say it:} a sentence of your own with \emph{los grandes almacenes}
  \item \emph{Recall:} say the Spanish for the street market, then the Spanish for the shopping centre, then say \emph{los grandes almacenes} again
\end{itemize}

\begin{tcolorbox}[breakable,colback=teal!4,colframe=teal!35!black,title={Before you move on}]
What does los grandes almacenes mean? (\textquotedblleft{}The department store\textquotedblright{}.) Say it once more.
\end{tcolorbox}
\section[festivo]{festivo — holiday (of a day: día festivo); festive}

\label{lesson:ES-C681-festivo}



\begin{quote}
Before the new one: say the Spanish for the shopping centre, then the Spanish for the department store.
\end{quote}

\subsection*{You'll want to know: festivo}
\textbf{festivo} — \textquotedblleft{}holiday (of a day: día festivo); festive\textquotedblright{}.

\begin{quote}
El museo cierra los días festivos.
\end{quote}

\begin{quote}
\emph{The museum is closed on public holidays.}
\end{quote}

On timetables: \emph{domingos y festivos} — Sundays and public holidays.

\begin{cousinweb}
\emph{festivo} is a learned borrowing from Latin \emph{fēstīvus} \textquotedblleft{}festive, merry\textquotedblright{} (Lewis \& Short), from \emph{fēstum} \textquotedblleft{}a holiday\textquotedblright{}.
\end{cousinweb}

\begin{grammarlens}[title={What you've built}]
One more word for telling, explaining and describing at B1.
\end{grammarlens}

\subsection*{Your turn}
\begin{itemize}
\raggedright
  \item \emph{Say it:} \emph{festivo}
  \item \emph{Say it:} \emph{festivo}, once more
  \item \emph{Say it:} a sentence of your own with \emph{festivo}
  \item \emph{Recall:} say the Spanish for the shopping centre, then the Spanish for the department store, then say \emph{festivo} again
\end{itemize}

\begin{tcolorbox}[breakable,colback=teal!4,colframe=teal!35!black,title={Before you move on}]
What does festivo mean? (\textquotedblleft{}Holiday (of a day: día festivo); festive\textquotedblright{}.) Say it once more.
\end{tcolorbox}
\section[atender]{atender — to serve (a customer); to attend to; to pay attention}

\label{lesson:ES-C681-atender}



\begin{quote}
Before the new one: say the Spanish for the department store, then the Spanish for holiday.
\end{quote}

\subsection*{You'll want to know: atender}
\textbf{atender} — \textquotedblleft{}to serve (a customer); to attend to; to pay attention\textquotedblright{}.

\begin{quote}
¿Ya le atienden?
\end{quote}

\begin{quote}
\emph{Are you being served?}
\end{quote}

The e becomes ie under stress: \emph{atiendo}, \emph{atiende}.

\begin{cousinweb}
\emph{atender} comes from Latin \emph{attendere} \textquotedblleft{}to stretch towards, to give heed\textquotedblright{} (EtymDB), built on \emph{tendere} \textquotedblleft{}to stretch\textquotedblright{}.
\end{cousinweb}

\begin{grammarlens}[title={What you've built}]
One more word for telling, explaining and describing at B1.
\end{grammarlens}

\subsection*{Your turn}
\begin{itemize}
\raggedright
  \item \emph{Say it:} \emph{atender}
  \item \emph{Say it:} \emph{atender}, once more
  \item \emph{Say it:} a sentence of your own with \emph{atender}
  \item \emph{Recall:} say the Spanish for the department store, then the Spanish for holiday, then say \emph{atender} again
\end{itemize}

\begin{tcolorbox}[breakable,colback=teal!4,colframe=teal!35!black,title={Before you move on}]
What does atender mean? (\textquotedblleft{}To serve (a customer); to attend to; to pay attention\textquotedblright{}.) Say it once more.
\end{tcolorbox}
